\subsection{Referências e incerteza das comparações}
\label{sec:protocolo_comparacao}

Nas linhas de treino dos codificadores, XGBoost e LightGBM \cite{chen2016xgboost,ke2017lightgbm} ajustam por mínimos quadrados um regressor por coordenada a \num{58} atributos agregados, como centroides com peso linear e quadrático, quantis, dispersões, assimetrias e canais mais claros. Cada regressor tem até \num{3000} árvores, com parada após \num{100} rodadas sem melhora na validação; os hiperparâmetros são os de menor erro 3-D médio de validação em \num{40} tentativas de TPE, \num{12} delas aleatórias \cite{akiba2019optuna}. Nos registros, a validação e o teste usaram luz com aumento de dados (perda e ativação espúria de canais, ruído, variação de ganho, afinamento de fótons), com \qty{4.1}{\percent} menos luz no teste. Este relatório reavalia os \textit{boosters} na luz limpa dos codificadores; com aumento, reproduz exatamente os valores registrados. A~busca e o treino, porém, usaram luz com aumento.

A rede da etapa anterior \cite{vicentin2024study} é citada, não reavaliada: usou outra partição (\qty{80}{\percent} e \qty{20}{\percent}), a geometria anterior à correção das tampas e, a julgar pelos gráficos, o referencial dos nomes de arquivo. O APEX \cite{marinho2025apex} usa outra simulação (\qty{10}{ppm} de xenônio, refletividade de \qty{79}{\percent} na gaiola de campo, luz reemitida pelos fotodetectores) e estimadores analíticos. Suas larguras de resolução nos próprios eixos, $\sigma_x$, $\sigma_y$ e $\sigma_z$, lidas da Figura~2 do artigo a \qty{0.26}{\cm} por \textit{pixel}, não têm definição publicada; a largura combinada $\sigma_{\mathrm{comb}}$ é a média quadrática delas. A comparação principal é em luz detectada igual, que não garante igual informação de posição: os \qty{530.1}{\pe}\ médios do teste equivalem a \qty{2.9}{\MeV} no APEX com o rendimento médio, \qty{180}{\pe\per\MeV}, e a \qty{4.9}{\MeV} com o mínimo, \qty{109}{\pe\per\MeV}, energia em que até a região mais escura do APEX dá essa luz, o que favorece o APEX; também se compara em energia igual, \qty{10}{\MeV}. Só $\sigma_{\mathrm{comb}}$ e $\sigma_z$ são comparados: a correspondência das larguras transversais com $x$ e $y$ não está estabelecida.

O baricentro $\bar{z}$, com calibração quadrática ajustada em \num{4000} eventos de treino sorteados, usa os mesmos \num{937} eventos de teste e o mesmo sorteio de fótons. Os intervalos de \qty{95}{\percent} vêm de reamostragem \textit{bootstrap} dos eventos de teste \cite{efron1979bootstrap}, pareada nos codificadores e \textit{boosters}, com \num{10000} reamostras comuns, e não pareada no baricentro, com \num{2000}; cobrem a amostragem do teste, não a variação entre treinos. Como os mesmos eventos de teste serviram à primeira comparação e ao diagnóstico que orientou a final (Seção~\ref{sec:dificuldades}), os erros de teste dos codificadores podem ser algo otimistas, o que os intervalos também não cobrem. Diferenças, razões e porcentagens são calculadas sem arredondamento.
